%%%%%%%%%%%%%%%%%%%%%%%%%%%%%%%%%%%%%%%%%%%%%%%%%%%%%%%%%%%%%%%%%%%%%%%%%%%%%%%
% *** Tabellen ***
%%%%%%%%%%%%%%%%%%%%%%%%%%%%%%%%%%%%%%%%%%%%%%%%%%%%%%%%%%%%%%%%%%%%%%%%%%%%%%%
%
\section{Tabellen}\label{sec:tabellen}
%
Dies ist eine einfache Gleichung

\begin{table}[h]
  \renewcommand{\arraystretch}{1.2} % kuenstliche Erhoehung der Zeilenhoehe
  \caption{Parameter - quasistatisches Vergleichsmodell mit Temperatureinfluss}
  \label{tab:einfache_tabelle}
  \vspace{0.5cm} 
  \begin{tabular}{|l|l|l|c|}
    \hline
                & \textbf{Parameter} & \textbf{Formelzeichen} & \textbf{Wert}\\
    \hline
    Lager         & Durchmesser (innen)         & $d_I$         & $100 mm$\\
                  & Durchmesser (außen)         & $d_a$         & $200 mm$\\
                  & Breite                      & $b$           & $200 mm$\\ 
                  & Kaltspiel ($20\grad C$)     & $\Delta r$    & $123 \mu m$\\
    \hline
    Ölzuführung   & Lage                        & $\varphi_{T}$     & $90\grad$\\
                  & Länge                       & $x_{T}$       & $15.708 mm \entspricht 18\grad$\\
                  & Breite (zentrisch)          & $y_{T}$       & $70 mm $\\
                  & Zuführdruck                 & $p_{T}$       & $0.07 bar$\\
                  & Zuführtemperatur            & $T_{T}$       & $40\grad C$\\
  \hline
  \end{tabular}
  \renewcommand{\arraystretch}{1.0} % kuenstliche Erhoehung der Zeilenhoehe zuruecknehmen
\end{table}
%
Die Referenz auf die Tabelle erfolgt durch \tab{tab:einfache_tabelle}.\\
Die Formatierung ist unter ./formatforlagen/shortcuts.tex bearbeitbar.\\\\\\\\
%
Dies ist eine komplexere Tabelle
\begin{table}[H]
    \renewcommand{\arraystretch}{1.5} % kuenstliche Erhoehung der Zeilenhoehe
    \caption{Zuordnung der Geschwindigkeits- und Beschleunigungskomponenten}%
    \label{tab:komplexere_tabelle}%
    \vspace{0.5cm}
    \begin{tabular}{|l|l|l|}
        \hline
        \multicolumn{3}{|c|}{Geschwindigkeits- und Beschleunigungskomponenten eines elastischen Körpers}\\%
        \hline %\hline %
        \multirow{2}{*}{Komponenten der Geschwindigkeit bzgl.} & $\un{r}_A' + \un{u}'$ & Translation\\%
        \cline{2-3}
                                                               & $\widetilde{\un{\un{\omega}}} \, (\un{r}_A + \un{c} + \un{u})$ & Rotation\\%
        \hline %\hline %
        \multirow{4}{*}{Komponenten der Beschleunigung bzgl.} & $\un{r}_A'' + \un{u}''$ & Translation\\%
        \cline{2-3}
                                                               & $\dot{\widetilde{\un{\un{\omega}}}} \, (\un{r}_A + \un{c} + \un{u})$ & Rotation\\%
        \cline{2-3}
                                                               & $2\,\widetilde{\un{\un{\omega}}} \, (\un{r}_A' + \un{u}')$ & gyroskopische Effekte\\%
        \cline{2-3}
                                                               & $\widetilde{\un{\un{\omega}}} \, \widetilde{\un{\un{\omega}}} \, (\un{r}_A + \un{c} + \un{u})$ & zentrifugale Effekte\\%
        \hline%\hline
        \multicolumn{1}{|l|}{Starrkörperbewegung: }& \multicolumn{2}{l|}{ $\un{r}_A + \un{c}$}\\%
        \multicolumn{1}{|l|}{elastische Deformation: } & \multicolumn{2}{l|}{ $\un{u}$}\\%
        \hline
    \end{tabular}
    \renewcommand{\arraystretch}{1.0} % kuenstliche Erhoehung der Zeilenhoehe zuruecknehmen
\end{table}
%
Die Referenz auf die Tabelle erfolgt durch \tab{tab:komplexere_tabelle}.\\
Die Formatierung ist unter ./formatforlagen/shortcuts.tex bearbeitbar.\\\\\\\\
%
\newpage
%
